\documentclass[10pt,a4paper]{amsart}
\usepackage[margin=19mm]{geometry}
\usepackage{amsmath,amssymb,mathtools}
\usepackage[scr=rsfs,cal=euler]{mathalfa}
\usepackage{xfrac}
\usepackage[unicode,hidelinks,bookmarks=false]{hyperref}
\newcommand{\ie}{i.e.\ }
\newcommand{\cf}{cf.\ }
\newcommand{\resp}{respectively\ }
\newcommand{\Subsection}{\subsection}
\providecommand{\nobreakdash}{\nobreak\hbox{-}}
\setlength{\emergencystretch}{2em}
\multlinegap0pt
\usepackage{amssymb}
\usepackage[shortlabels]{enumitem}
\usepackage{tikz}
\usepackage[normalem]{ulem}
\newtheorem{lem}{Lemma}
\newcommand{\lt}{\left}
\newcommand{\na}{\nabla}
\newcommand{\pl}{\partial}
\newcommand{\rt}{\right}
\title{Quasi-Spherical Metrics and the Static Minkowski Inequality}
\author{Brian Harvie, Ye-Kai Wang}
\date{Source: arXiv:2403.06216v3 (CC BY 4.0)}
\begin{document}
\maketitle

\noindent\emph{Source and licence.} Quasi-Spherical Metrics and the Static Minkowski Inequality. Version arXiv:2403.06216v3 (CC BY 4.0), distributed by arXiv under the Creative Commons Attribution 4.0 International licence, http://creativecommons.org/licenses/by/4.0/, as recorded in the arXiv licence metadata for this exact version and shown on the abstract page https://arxiv.org/abs/2403.06216v3. Full licence notice: \texttt{licenses/arxiv-2403.06216-CC-BY-4.0.txt}.\par\smallskip

\noindent\textit{Editorial scope.} Counted unit: the article's lem inequality (source0.tex lines 446--453). The article's complete original proof of it is delivered with it (source0.tex lines 455--539, one \texttt{proof} environment of 85 lines). No mathematics is written, repaired or re-derived: the statement and the proof are the article's own lines, in the article's own notation, and no retained line is substituted or re-flowed.  Retained uncounted as the source-local context the proof needs: lines 407--409; lines 422--424.  Nothing was omitted from any retained span, so the package carries no omission record.  No retained line contains a citation, so the wrapper carries no bibliography and no bibliography entry.  No retained line contains a figure environment, an image-inclusion command or drawing code, so no image is delivered and the package declares no asset dependency.  Editorial formatting only, in addition: an AMS \texttt{amsart} preamble that restates the article's own declarations and macros used by the retained lines, namely the environments \texttt{lem} on separate counters, together with the standard packages those lines need.  The retained environments print in this document as Lemma 1 in order of appearance, whereas the article prints the same items with the numbering of its own full document.  The complete native source of the article as distributed in the arXiv e-print for this exact version is bundled unchanged as \texttt{sources/155855/source0.tex}.
\begin{align}\label{C0C1bound}
r_0 e^{\frac{t}{n-1}} \le r \le (r_0 + r_1(t)) e^{\frac{t}{n-1}}, \quad \langle \pl_r ,\nu \rangle > \eta(t)
\end{align}

Let $l > 0$. For $f = r^{-l}$ and $X = r \frac{\pl}{\pl r}$, we have \begin{align}\label{D2f}
D_iD_j f -(\Delta_g f) g_{ij}  = l \lt[ (l+2) r^{-l-4} x^ix^j + (n-3-l) r^{-l-2}\delta_{ij} \rt] + o(r^{-l-2})
\end{align}

\begin{lem}\label{differential inequality}
Let $u = H^{-1} f^{-1} \langle X,\nu \rangle^{-2}$. Then there exists a large time $T_2$ and small $l$, depending only on $\epsilon(r)$ and $n$, such that 
\begin{align}
\pl_t u &< \frac{1}{H^2}\Delta u - \frac{2}{H^2} u^{-1} |\na u|^2 - \frac{2}{H^3} \na^a H \na_a u \notag\\
&= \na^a \lt( \frac{1}{H^2} \na_a u \rt) - \frac{2}{H^2} u^{-1}|\na u|^2
\end{align}
on $\Sigma_t, t \ge T_2$. 
\end{lem}

\begin{proof}
Note that the curvatures of $g$ satisfies $R_{ijkl} = o(r^{-2})$. We compute
\begin{align*}
\pl_t \langle X,\nu \rangle = \frac{1}{H} \lt( 1 + \varepsilon(\nu,\nu) \rt) - \langle X, \na (\frac{1}{H}) \rangle,
\end{align*}
\begin{align}\label{gradXnu}
\pl_b \langle X,\nu \rangle = \varepsilon(\pl_b, \nu) + h_b^c \langle X,\pl_c \rangle, 
\end{align}
and by Codazzi equation
\begin{align*}
\Delta \langle X,\nu \rangle &= \sigma^{ab} \lt[ \pl_a \lt( \varepsilon(\pl_b, \nu) + h_b^c \langle X,\pl_c \rangle \rt) - \Gamma_{ab}^c \lt( \varepsilon(\pl_c,\nu) + h_c^d \langle X, \pl_d \rangle \rt)\rt]\\
&= \sigma^{ab} D_{\pl_a} \varepsilon (\pl_b,\nu) - H \varepsilon(\nu,\nu) + h^{bc} \varepsilon(\pl_b,\pl_c) \\
&\quad + \langle X,\na H \rangle + R(X^T, \nu) + H + h^{ac} \varepsilon(\pl_a, \pl_c) - |A|^2 \langle X,\nu \rangle.
\end{align*}
Here $X^T$ denotes the tangential component of $X$ along $\Sigma_t$.  As a result, the evolution equation of $\langle X,\nu \rangle$ is 
\begin{align*}
\pl_t \langle X,\nu \rangle = \frac{1}{H^2} \lt( \Delta \langle X,\nu \rangle + |A|^2 \langle X,\nu \rangle - R(X^T,\nu) + A * \varepsilon + o(r^{-1}) \rt).
\end{align*}
Using the relation of Laplacians on a hypersurface, we get 
\begin{align*}
\pl_t f = \frac{2}{H} \frac{\pl f}{\pl\nu} + \frac{1}{H^2} \lt( \Delta f - \Delta_g f + D^2 f(\nu,\nu) \rt).
\end{align*}

We are ready to compute the evolution equation of $H f \langle X,\nu \rangle^2$.
\begin{align*}
&\pl_t \lt( H f \langle X,\nu \rangle^2 \rt) \\
&= \lt( \frac{1}{H^2}\Delta H - \frac{2}{H^3} |\na H|^2 - \frac{|A|^2}{H} - \frac{R(\nu,\nu)}{H} \rt) f \langle X,\nu\rangle^2 \\
&\quad + \frac{1}{H^2} \lt(\Delta f +  2H \frac{\pl f}{\pl\nu} + D^2 f(\nu,\nu) - \Delta_g f \rt) \\
&\quad + 2 H f \langle X,\nu \rangle \cdot \frac{1}{H^2} \lt( \Delta \langle X,\nu \rangle + |A|^2 \langle X,\nu \rangle - R(X^T,\nu) + A * \varepsilon + o(r^{-1}) \rt)\\
&= \frac{1}{H^2} \lt[ \Delta (H f \langle X,\nu \rangle^2) - 2 \na H \cdot \na f \langle X,\nu \rangle^2 - 2 H \na f \cdot \na \langle X,\nu \rangle^2 - 2 f \na H \cdot \na\langle X,\nu \rangle^2 \rt] \\
&\quad - \frac{2}{H^3} |\na H|^2 \cdot f \langle X,\nu \rangle^2 - \frac{f}{H} |\na \langle X,\nu \rangle|^2  \\
&\quad + \frac{\langle X,\nu \rangle}{H} \Big[  f |A|^2 \langle X,\nu \rangle - f R(\nu,\nu) \langle X,\nu \rangle + 2 H \frac{\pl f}{\pl\nu} \langle X,\nu \rangle \\
&\qquad\qquad\qquad\qquad + (D^2 f(\nu,\nu) - \Delta_g f)\langle X,\nu \rangle - 2f R(X^T,\nu) + 2f A * \varepsilon + f \cdot o(r^{-1}) \Big].
\end{align*}
Note that \eqref{C0C1bound} implies 
\begin{align*}
|X^T| = \langle X,\nu \rangle \cdot o(1)
\end{align*} Hence by \eqref{gradXnu} and Cauchy-Schwarz, we have \begin{align*}
|\na\langle X,\nu \rangle|^2  \le 2 |A^2| |X^T|^2 + o(1) = |A|^2 \langle X,\nu \rangle^2 \cdot o(1) + o(1).
\end{align*}
Moreover, note that $R(\nu,\nu) = o(r^{-2})$ and $R(X^T,\nu) = o(r^{-1})$. We thus obtain
\begin{align*}
\pl_t \lt( Hf\langle X,\nu \rangle^2 \rt) &= \frac{1}{H^2}\Delta (Hf\langle X,\nu \rangle^2) - \frac{2}{H^3} \na H \cdot \na (Hf\langle X,\nu \rangle^2) + \frac{\langle X,\nu\rangle}{H} \cdot I
\end{align*}
where
\begin{align*}
I &= -4 \na f \cdot \na\langle X,\nu \rangle + f |A|^2 \langle X,\nu \rangle (1 - o(1)) + 2H \frac{\pl f}{\pl\nu} \langle X,\nu \rangle \\
&\quad +(D^2 f(\nu,\nu) - \Delta_g f)\langle X,\nu \rangle + 2f A * \varepsilon + f \cdot o(r^{-1})
\end{align*}

We claim that $I > 0$ if $t$ is sufficiently large and $l$ is sufficiently small. First of all, we have 
\begin{align*}
|\na f| &= |Df| \cdot o(1) = f \cdot o(r^{-1}) \\
\frac{\pl f}{\pl\nu} &= - |Df| (1 + o(1)) = - lf r^{-1} (1 + o(1))
\end{align*}
as a consequence of \eqref{C0C1bound}.
By \eqref{gradXnu} and Cauchy-Schwarz, we have
\begin{align*}
- \na f \cdot \na \langle X,\nu \rangle &= - A(X^T,\na f) - \varepsilon (\na f,\nu) \ge  - |A|\langle X,\nu \rangle f \cdot o(r^{-1}) - f \cdot o(r^{-1}) \\
&\ge -\frac{1}{16} f |A|^2 \langle X,\nu \rangle  - f \cdot o(r^{-1}).
\end{align*}
By Cauchy-Schwarz, we have
\begin{align*}
2H \frac{\pl f}{\pl\nu} \langle X,\nu \rangle &\ge - \frac{H^2}{4(n-1)} f \langle X,\nu \rangle - 4(n-1) \lt( \frac{\pl f}{\pl\nu}\rt)^2 \frac{\langle X,\nu \rangle}{f} \\
&\ge - \frac{1}{4} f |A|^2 \langle X,\nu \rangle - 4(n-1)l^2 \frac{f}{r^2} \langle X,\nu \rangle (1 + o(1))
\end{align*}
and
\begin{align*}
2f A * \varepsilon \ge - \frac{1}{4}f |A|^2 \langle X,\nu \rangle + f \langle X,\nu \rangle^{-1} \cdot o(1)
\end{align*}
Putting these together, we get
\begin{align*}
I \ge f|A|^2 \langle X,\nu \rangle \lt( \frac{1}{4} - o(1) \rt) + \lt( D^2 f(\nu,\nu) - \Delta_g f - 4(n-1) \frac{l^2}{r^2} f - f \cdot o(r^{-2}) \rt)\langle X,\nu \rangle.
\end{align*}
By \eqref{D2f}, $D^2 f(\nu,\nu) - \Delta_g f = l(n-1) r^{-l-2} + o(r^{-l-2})$ and its follows that the second term on the right-hand side is positive if $l$ is small enough and $r$ is large enough (equivalently $t$ large enough). This completes the proof of the claim and we arrive at
\begin{align*}
\pl_t \lt( Hf\langle X,\nu \rangle^2 \rt) &> \frac{1}{H^2}\Delta (Hf\langle X,\nu \rangle^2) - \frac{2}{H^3} \na H \cdot \na (Hf\langle X,\nu \rangle^2).
\end{align*}

For $u = \lt( Hf\langle X,\nu \rangle^2\rt)^{-1}$, we have
\begin{align*}
\pl_t u &< - (Hf \langle X,\nu \rangle^2)^{-2} \lt( \frac{1}{H^2}\Delta (Hf\langle X,\nu \rangle^2) - \frac{2}{H^3} \na H \cdot \na (Hf\langle X,\nu \rangle^2) \rt) \\
&= \frac{1}{H^2}\Delta u - \frac{2}{H^2} u^{-1} |\na u|^2 - \frac{2}{H^3} \na H \cdot \na u
\end{align*}
\end{proof}

\end{document}
